\documentclass[11pt]{article}
\usepackage[margin=1in]{geometry}
\usepackage{amsmath,amssymb,amsthm,mathtools,microtype}
\usepackage[hidelinks]{hyperref}
\hypersetup{pdftitle={The outer center-pedal area product in elliptic billiards},pdfauthor={Alec Kriebel}}
\newtheorem{theorem}{Theorem}
\newtheorem{lemma}[theorem]{Lemma}
\newcommand{\sn}{\operatorname{sn}}
\newcommand{\cn}{\operatorname{cn}}
\newcommand{\dn}{\operatorname{dn}}
\newcommand{\cs}{\operatorname{cs}}
\newcommand{\Area}{\mathcal A}
\newcommand{\ii}{\mathrm i}
\title{The outer center-pedal area product\\in elliptic billiards}
\author{Alec Kriebel\\Independent researcher\\\small\href{https://orcid.org/0009-0001-9320-500X}{ORCID: 0009-0001-9320-500X}}
\date{16 October 2026}
\begin{document}
\maketitle
\begin{abstract}
For a periodic billiard family between strictly nested confocal ellipses,
form the outer tangent polygon and project the ellipse center onto its side
lines. We prove constancy of the product of their signed areas when the
period is not divisible by four, including coprime stars and repeated
traversals. This proves the assertion labelled $k_{303,b}$ in the 2020
invariant list and its 2021 published edition. The actual center-foot map
is a fixed linear image of the edge-vector polygon of a phase-shifted
billiard. Explicit area formulas reduce the assertion to a classical odd
cyclic Jacobi identity, also giving the invariant's value. The same full
theorem is stated and proved in Ferudun's contemporary version 1.1; no
novelty, historical independence or priority claim is made here. A source
and verification package accompanies this unrefereed, extensively AI-assisted
proof and verification note.
\end{abstract}

\section{The precise invariant}
For an ordered closed polygon $W=(W_0,\ldots,W_{N-1})$ define
\begin{equation}\label{eq:area}
 \Area(W)=\frac12\sum_{j=0}^{N-1}\det(W_j,W_{j+1}),\qquad W_N=W_0.
\end{equation}
Thus stars retain trajectory order and signed shoelace area. Let
\[
 E:\frac{x^2}{a^2}+\frac{y^2}{b^2}=1,\quad a>b>0,\qquad
 C:\frac{x^2}{\alpha^2}+\frac{y^2}{\beta^2}=1,
 \quad a^2-\alpha^2=b^2-\beta^2=\lambda\in(0,b^2).
\]
For a billiard orbit $P_j$ on $E$ tangent to $C$, let $\ell_j$ be the
tangent to $E$ at $P_j$, let $R_j=\ell_j\cap\ell_{j+1}$, and let $Q_j$
be the perpendicular foot of $O=(0,0)$ on $\ell_j$. The line $\ell_j$
is the side line $R_{j-1}R_j$ of the outer polygon. Write
$A'=\Area(R)$ and $A'_O=\Area(Q)$.

\begin{theorem}\label{thm:main}
Suppose the fixed pair $(E,C)$ supports a periodic billiard family of least
period $N\geq3$, with $N\not\equiv0\pmod4$. Then $A'A'_O$ is independent
of the initial vertex. The assertion includes every admissible coprime star
winding. It also holds for a listed traversal length not divisible by four
that repeats a shorter orbit, with the product multiplied by the square
of the repetition count.
\end{theorem}

This is $k_{303,b}$ of Reznik--Garcia--Koiller
\cite[Table 4, p.~6]{RGK} and \cite[Table 4, p.~346]{RGKfinal};
the construction and signed-area convention are in their Sections 2 and 3.4.
The conjecture first appeared in their 2020 arXiv source; the companion
journal edition appeared in 2021. The source concerns a confocal pair of
ellipses. Our theorem excludes hyperbolic, focal and collapsed caustics.
Feet are on the full boundary tangent lines, not on the original billiard
chords. Areas of spatially sorted vertices or filled star lobes are different
quantities.

The exact full theorem, with these parity classes, is also proved in
Ferudun \cite[Theorem 2 and Section 6]{Ferudun}, version 1.1. That version,
DOI \href{https://doi.org/10.5281/zenodo.23088516}{10.5281/zenodo.23088516},
is distinguished from version 1.0, which gives the primitive even
arbitrary-point result and hence its centered even subset. Section 4 records
the scope of the priority evidence. Here we give the explicit edge-vector
and classical-identity reduction, rather than assert a first resolution.

\section{Area formulas in canonical coordinates}
All Jacobi functions below have modulus $k$, not parameter $k^2$. Set
\[
 k=\frac{\sqrt{\alpha^2-\beta^2}}{\alpha},\quad
 k'=\frac{\beta}{\alpha},\quad K=K(k),\quad K'=K(k').
\]
Stachel's canonical parametrization
\cite[Theorem 4.3, equation (4.9), p.~1614]{Stachel} gives, after choosing
orientation,
\begin{gather}\label{eq:canonical}
 P_j=P(u+2jv),\quad P(z)=(-a\sn z,b\cn z),\\
 v=\frac{2\tau K}{N}\in(0,K),\quad \gcd(\tau,N)=1,\quad
 0<\tau<N/2,\qquad
 a=\frac{\alpha\dn v}{\cn v},\quad b=\frac{\beta}{\cn v}.\notag
\end{gather}
Only windings supported by the fixed pair are asserted. Put
\[
 f(t)=\frac{\sn t\cn t}{\dn t},\qquad
 T(z)=\sum_{j=0}^{N-1}\dn(z+2jv).
\]
Cyclic reindexing gives $T(z+2v)=T(z)$. Jacobi addition
\cite[Section 22.8]{DLMF} gives, for $H(z)=(-A\sn z,B\cn z)$,
\begin{equation}\label{eq:det}
 \det(H(z-t),H(z+t))
 =\frac{2AB\sn t\cn t\dn z}{1-k^2\sn^2z\sn^2t}
 =ABf(t)\bigl(\dn(z-t)+\dn(z+t)\bigr).
\end{equation}
Summing consecutive determinants proves $\Area(P)=abf(v)T(u)$.
For any indexed polygon, determinant expansion also gives
\[
 \Area((P_{j+1}-P_j)_j)
 =2\Area(P)-\frac12\sum_j\det(P_j,P_{j+2}).
\]
Applying \eqref{eq:det} with $t=2v$ and reindexing the sums yields
\begin{equation}\label{eq:edge}
 \Area((P_{j+1}-P_j)_j)=ab\bigl(2f(v)-f(2v)\bigr)T(u).
\end{equation}
This identity does not require stepping by two to form a single cycle.

The side $P_jP_{j+1}$ is tangent to $C$ at
$B(u+(2j+1)v)$, where $B(z)=(-\alpha\sn z,\beta\cn z)$.
Indeed both $P(z\pm v)$ satisfy
$-\sn(z)X/\alpha+\cn(z)Y/\beta=1$ by addition formulas.
Its pole with respect to $E$ is $R_j$, so elementary conic polarity gives
\[
 R_j=\operatorname{diag}(a^2/\alpha^2,b^2/\beta^2)B(u+(2j+1)v)
 =\operatorname{diag}(a/\alpha,b/\beta)P(u+(2j+1)v).
\]
Taking the fixed linear area scale,
\begin{equation}\label{eq:outer}
 A'(u)=\frac{a^2b^2}{\alpha\beta}f(v)T(u+v).
\end{equation}
The intersections $R_j$ are finite: parallel endpoint tangents would force
antipodal endpoints, whose center-crossing chord cannot touch the strictly
interior ellipse $C$.

\subsection*{The actual perpendicular foot}
The tangent at $P(z)$ has equation $n(z)\cdot X=1$, where
$n(z)=(-\sn z/a,\cn z/b)$. Thus
\begin{equation}\label{eq:foot}
 Q(z)=\frac{n(z)}{n(z)\cdot n(z)}
 =\frac{(-ab^2\sn z,a^2b\cn z)}{a^2-(a^2-b^2)\sn^2z}.
\end{equation}
Set $w=K-v\in(0,K)$. Quarter-period formulas give
\[
 \sn w=\frac{\cn v}{\dn v},\qquad
 \cn w=\frac{k'\sn v}{\dn v}>0,\qquad \dn w=\frac ba.
\]
The sum formulas for $P(z+w)+P(z-w)$ have denominator
\[
 D(z)=1-k^2\sn^2z\sn^2w
 =\frac{a^2-(a^2-b^2)\sn^2z}{a^2}
\]
and numerator $(-2a\sn z\cn w\dn w,2b\cn z\cn w)$.
Using $P(z-2K)=-P(z)$ now proves the pointwise identity
\begin{equation}\label{eq:edgefoot}
 Q(z)=L\bigl(P(z+K-v)-P(z+K+v)\bigr),\quad
 L=\frac{\operatorname{diag}(b/a,1)}{2\cn w},\quad
 \det L=\frac{b}{4a\cn^2w}.
\end{equation}
The global minus sign relative to the edge vector preserves area in two
dimensions. Therefore \eqref{eq:edge} proves
\begin{equation}\label{eq:pedal}
 A'_O(u)=\frac{b^2}{4\cn^2(K-v)}
 \bigl(2f(v)-f(2v)\bigr)T(u+K-v).
\end{equation}
All real denominators used above are positive or nonzero. No polygon area
has been divided out.

The underlying geometry is classical. With $e_j$ the forward unit chord
velocity and $J=\sqrt\lambda/(ab)$ the Joachimsthal integral, reflection
gives $Q_j=(e_{j-1}-e_j)/(2J)$, as follows directly from
\cite[proof of Theorem 4]{AST} and in the center-foot telescope of
\cite[Theorem 4.1]{BT}. The skew hodograph
$\operatorname{diag}(a,b)e_j$ is described in
\cite[Lemma 3.1 and Remark 3.2]{AST}. In the present coordinates its exact
phase is $P(u+(2j+1)v+K)$. Equation \eqref{eq:edgefoot} verifies this
mechanism directly for the same caustic and fixes its normalization.

\section{The classical trace identity and proof of the theorem}
\begin{lemma}\label{lem:trace}
Let $q$ be positive and odd, and $0<k<1$. Define
\[
 U_q(z)=\sum_{j=0}^{q-1}\dn(z+2Kj/q),\qquad
 C_q=\prod_{j=1}^{(q-1)/2}\cs^2(2Kj/q),\quad \cs=\cn/\sn.
\]
The empty product is one. Then, as a meromorphic identity,
\begin{equation}\label{eq:trace}
 U_q(z)U_q(z+K)=\frac{(k')^q}{C_q^2}.
\end{equation}
\end{lemma}
\begin{proof}
Khare--Lakshminarayan--Sukhatme
\cite[Appendix A, equation (86) in arXiv v1; (85) in v2]{KLS}
give the odd cyclic product-to-sum identity
\[
 \prod_{j=0}^{q-1}\dn(z+2Kj/q)=C_q U_q(z).
\]
Multiply it by its $K$-translate and use
$\dn(t+K)=k'/\dn t$. Meromorphic continuation handles isolated zeros
and poles. For $q=1$ this is just the quarter-period formula.
\end{proof}
This trace relation is prior mathematics, also connected to the generalized
Landen formulas in \cite{KS}. For completeness its product-to-sum precursor
can be checked directly. The product $V(z)$ and sum $U_q(z)$ are
$2K$-periodic and $2\ii K'$-antiperiodic. Their poles are simple at the
shifts of $\ii K'$ by the real subgroup $H=\{2Kj/q\}$ modulo $2K$.
At $\ii K'$, the standard formulas
$\operatorname{Res}\dn=-\ii$ and
$\dn(\ii K'+t)=-\ii\cs t$ give
\[
 \frac{\operatorname{Res}_{\ii K'}V}{\operatorname{Res}_{\ii K'}U_q}
 =(-\ii)^{q-1}\prod_{j=1}^{q-1}\cs(2Kj/q)=C_q,
\]
since $\cs(2K-t)=-\cs t$. The same ratio holds at every pole by subgroup
periodicity and imaginary antiperiodicity. Thus $V-C_qU_q$ is entire
elliptic, hence constant, and its antiperiodicity forces zero. This is a
reproof of the classical identity, not a new special-function assertion.
The meromorphic pole-cancellation method itself is established, for example
in \cite{KLS,AST,CC}.

\begin{proof}[Proof of Theorem \ref{thm:main}]
First consider a primitive orbit and put
$q=N/\gcd(N,2)$, $r=N/q$. Coprimality in \eqref{eq:canonical} means that
the shifts $2jv$ modulo $2K$ exhaust $H$ exactly $r$ times. Consequently
$T(z)=rU_q(z)$. The hypothesis $N\not\equiv0\pmod4$ says precisely that
$q$ is odd. Because $T$ has period $2v$,
$T(u+K-v)=T(u+K+v)$. Multiplying \eqref{eq:outer} and
\eqref{eq:pedal}, then using Lemma \ref{lem:trace}, gives the explicit value
\begin{equation}\label{eq:constant}
 A'A'_O=
 \frac{a^2b^4 f(v)\bigl(2f(v)-f(2v)\bigr)}
 {4\alpha\beta\cn^2(K-v)}\,
 \frac{r^2(k')^q}{C_q^2},
\end{equation}
independent of $u$. These determinant and subgroup computations apply to
stars without reordering. Reversing traversal changes both area signs and
preserves their product. A $d$-fold traversal repeats each shoelace sum $d$
times. If its listed length is not divisible by four, neither is its least
period, so the primitive theorem applies and its product is multiplied by
$d^2$.
\end{proof}

The coefficient in \eqref{eq:constant} has units length to the fourth
power, as required. Nondegenerate one- or two-bounce orbits are impossible
here: the two-bounce segment passes through $O$ and cannot touch $C$.
The circular limit is separately elementary: all fixed-step regular-star
families differ by rigid rotation, so both signed areas are constant.
There is no assertion for a general period divisible by four. For example,
on $x^2/4+y^2=1$ the diamond with vertices $(2,0),(0,1),(-2,0),(0,-1)$
and the rectangle with vertices $(\pm4/\sqrt5,\pm1/\sqrt5)$ in boundary
order are billiards with the same caustic $\lambda=4/5$.
Their outer/pedal products are respectively $8\cdot4=32$ and
$10\cdot5=50$. Thus the omitted parity cannot simply be dropped.

\section{Attribution, evidence and limitations}
The note proves the originally conjectured invariant in the stated domain.
It does not establish that this is a new theorem or the first proof.
Ferudun's version 1.1 \cite{Ferudun} contains the exact full result.
Its Zenodo record was created on 1 October 2026 at 18:12:36 UTC; version 1.0
was created at 16:56:14 UTC. GitHub server check-suite timestamps place
receipt of this note's unchanged candidate proof tree at 05:35:19 UTC that
day; PR 209 was created at 05:48:56 UTC. These observations do not certify
past public repository visibility, exclude an earlier website posting, or
establish historical independent non-derivation. The bounded priority audit
and its byte-bound chronology are included in the supplement. They leave
historical ordering and novelty undetermined. The old trace identities,
foot/velocity and hodograph mechanisms, canonical coordinates and
meromorphic method are expressly credited above. Chavez-Caliz's even
original/outer area product \cite[Theorem 6]{CC} concerns a different pair
of areas; its auxiliary $2q$-family also yields trace constancy, as detailed
in the supplement.

The verification package preserves the original candidate and inherited
review unchanged, three fresh mathematical/source audits and two priority
audits, derivations, original and new exact controls, replay receipts and
failure corrections. Exact symbolic calculations and arbitrary-precision
samples support the analytic proof; finite samples are not universal proofs
or interval-certified computations. In particular the period-four examples
check the actual reflection and common-caustic hypotheses. The package
states software versions and how to rerun the mathematical checks without
network access; source retrieval and release operations are separate from
those checks. Third-party primary papers and local software environments
are not redistributed.

\paragraph{AI use and review status.}
AI tools were used extensively for discovery, derivation, programming,
literature searches, writing and adversarial checking, including separately
prompted automated reviewers. This is an unrefereed preprint. No independent
human peer review or formal proof verification is claimed. Automated reviews
and exact controls do not constitute either. Alec Kriebel is the named author;
the preserved audit records state the actual scope of the checks.

\begin{thebibliography}{99}
\bibitem{RGK} D. Reznik, R. Garcia and J. Koiller,
\emph{Eighty New Invariants of N-Periodics in the Elliptic Billiard},
arXiv:2004.12497v11 (2020),
\url{https://arxiv.org/abs/2004.12497v11}.
\bibitem{RGKfinal} D. Reznik, R. Garcia and J. Koiller,
\emph{Fifty New Invariants of N-Periodics in the Elliptic Billiard},
Arnold Math. J. \textbf{7} (2021), 341--355,
\href{https://doi.org/10.1007/s40598-021-00174-y}{doi:10.1007/s40598-021-00174-y}.
\bibitem{Stachel} H. Stachel, \emph{On the motion of billiards in ellipses},
European J. Math. \textbf{8} (2022), 1602--1622,
\href{https://doi.org/10.1007/s40879-021-00524-2}{doi:10.1007/s40879-021-00524-2}.
\bibitem{KLS} A. Khare, A. Lakshminarayan and U. Sukhatme,
\emph{Cyclic Identities Involving Jacobi Elliptic Functions. II},
J. Math. Phys. \textbf{44} (2003), 1822--1841;
arXiv:math-ph/0207019 (2002),
\href{https://doi.org/10.1063/1.1560856}{doi:10.1063/1.1560856}.
\bibitem{KS} A. Khare and U. Sukhatme,
\emph{Connecting Jacobi elliptic functions with different modulus parameters},
Pramana \textbf{63} (2004), 921--936.\\
Author version: \emph{Generalized Landen Transformation Formulas for Jacobi Elliptic Functions}, arXiv:math-ph/0312074v1 (2003),
\href{https://doi.org/10.1007/BF02704331}{doi:10.1007/BF02704331}.
\bibitem{AST} A. Akopyan, R. Schwartz and S. Tabachnikov,
\emph{Billiards in ellipses revisited},\\
arXiv:2001.02934v2 (2020),
\href{https://doi.org/10.1007/s40879-020-00426-9}{doi:10.1007/s40879-020-00426-9}.
\bibitem{BT} M. Bialy and S. Tabachnikov,
\emph{Dan Reznik's identities and more}, arXiv:2001.08469v2 (2020),
\href{https://doi.org/10.1007/s40879-020-00428-7}{doi:10.1007/s40879-020-00428-7}.
\bibitem{CC} A. Ch\'avez-Caliz,
\emph{More About Areas and Centers of Poncelet Polygons},
Arnold Math. J. \textbf{7} (2021), 91--106,
\href{https://doi.org/10.1007/s40598-020-00154-8}{doi:10.1007/s40598-020-00154-8}.
\bibitem{Ferudun} A. Ferudun,
\emph{Outer-Pedal Area Products in Elliptic Billiards}, version 1.1 (2026),
Theorem 2 and Section 6,
\href{https://doi.org/10.5281/zenodo.23088516}{doi:10.5281/zenodo.23088516}
(concept DOI 10.5281/zenodo.23087220).
\bibitem{DLMF} NIST Digital Library of Mathematical Functions,
Sections 22.4, 22.6 and 22.8, \url{https://dlmf.nist.gov/22}.
\end{thebibliography}
\end{document}
